\chapter{Mapping genetic architecture of methylation quantitative trait loci}

\section{Genetic architecture of epigenetic regulation}

Genetic variants exert profound regulatory influences on local and distal DNA methylation levels. Previous comprehensive mapping initiatives \cite{Smith2024} demonstrated that methylation quantitative trait loci (meQTL) are enriched in disease-associated regions. Linking these variants to aging markers \cite{Horvath2013} establishes mechanistic pathways connecting sequence variation to epigenetic age acceleration. Complementary approaches are discussed in \cite{Cell2020}.

\subsection{Ancestry-stratified fine-mapping procedures}

We implemented Bayesian fine-mapping algorithms to dissect causal variant credible sets across ancestrally distinct cohorts. Consistent with \cite{Pubmed2024}, trans-ancestry fine-mapping substantially narrowed the credible set sizes by exploiting differential linkage disequilibrium architecture.

\begin{quotation}
\begin{singlespace}
  Deciphering the causal interplay between sequence-level genetic variation and dynamic epigenetic modifications is foundational to interpreting GWAS discoveries and constructing biologically robust, multi-ancestry risk prediction models.
\end{singlespace}
\end{quotation}

\subsection{Colocalization and functional validation}

Colocalization analyses confirmed shared genetic etiology between meQTL and complex traits in over 42\% of tested genomic loci.
